\documentclass[12pt]{article}

\usepackage[left=1.25in,right=1.25in,top=1.2in,bottom=1.2in]{geometry}
\usepackage{setspace}
\usepackage{microtype}
\usepackage[T1]{fontenc}
\usepackage[utf8]{inputenc}
\usepackage{times}
\usepackage{amsmath,amssymb}
\usepackage{csquotes}
\usepackage{hyperref}

\setstretch{1.2}
\setlength{\parindent}{0pt}
\setlength{\parskip}{6pt}

\title{\textbf{On Gaudiness}\\
\large Taste as Expressive Restraint}

\author{Flyxion}
\date{\today }

\begin{document}
\maketitle

\begin{abstract}
Contemporary discussions of bad taste frequently attribute gaudiness, tackiness, or aesthetic excess to technological causes, particularly the rise of generative artificial intelligence. This essay argues that such diagnoses mistake the symptom for the cause. Gaudiness is not a function of energy expenditure, computational excess, or novelty of medium, but rather of failed restraint, insufficient socialization, and an inability to suppress expressive impulses that are locally intelligible yet globally inappropriate. Drawing analogies from animal behavior, historical portraiture, information design, and AI-generated audio narration, this essay develops a theory of tackiness as a failure of modulation rather than production. Taste, on this account, is defined negatively: as the disciplined capacity to decide what not to express, even when expression is technically possible or locally justified.
\end{abstract}

\newpage
\section{Introduction: Against Technological Explanations of Bad Taste}

The terms \emph{tacky} and \emph{gaudy} occupy a curious position in aesthetic discourse. They are widely used, immediately understood, and yet rarely analyzed with any conceptual precision. When analysis does occur, it is often displaced onto external causes: mass culture, commercialization, algorithmic optimization, or, more recently, artificial intelligence. In these accounts, aesthetic failure is treated as a consequence of scale, speed, or automation, as though taste were a fragile substance dissolved by too much throughput.

This essay rejects that framing. Gaudiness is not caused by abundance, nor by novelty of medium. It is not a thermodynamic problem, nor a computational one. The presence of excess detail, ornament, or expressive force is not itself an aesthetic failure. What distinguishes gaudiness from richness is not quantity but regulation. The central claim of this essay is that tackiness emerges when expressive systems lack adequate mechanisms of inhibition, calibration, and contextual awareness. The failure is not that too much is produced, but that too little is withheld.

To make this argument concrete, I begin not with art, but with behavior. Aesthetic judgment, like social judgment, is best understood as a form of trained restraint. Before turning to images, diagrams, or AI-generated media, it is instructive to consider cases in which expressive enthusiasm is plainly visible yet intuitively unappealing, even when it is benign in intent. Such cases clarify that the discomfort provoked by gaudiness is not moral, nor energetic, but structural.

\section{Impulse Without Modulation: A Behavioral Analogy}

On a daily walk to get coffee, I pass a neighboring house with a dog who reliably launches itself at the window, barking and scratching with maximal intensity. When encountered in person, the behavior is similar: loud, insistent, and poorly regulated. There is little reason to believe the animal is aggressive. It likely recognizes me, understands that I pose no threat, and may even associate my presence with novelty, attention, or the possibility of a treat. The display appears to be a greeting rather than a warning.

Nevertheless, the behavior is unsettling. I generally ignore the dog, not out of fear or irritation, but because engagement would escalate the situation. Responding would intensify the pulling on the leash, provoke embarrassment or scolding from the owner, and amplify the very spectacle that makes the encounter uncomfortable. The problem is not the dog’s friendliness. It is the absence of restraint in how that friendliness is expressed.

A similar pattern appears in another small dog I know, which habitually mouths at hands without intent to injure and periodically bolts into traffic or unfenced yards. The animal is affectionate and curious, but its lack of impulse control renders it dangerous to itself and potentially alarming to strangers. A baby, an unfamiliar adult, or a startled passerby could easily misinterpret the behavior, leading to harm despite the absence of malice.

These cases illustrate a general principle: expressive energy is not inherently objectionable. What produces discomfort is expressive energy that has not been socialized into a shared interpretive framework. The dogs are not bad. They are untrained. Their behavior is locally intelligible but globally disruptive. It is legible as enthusiasm, yet illegible as etiquette.

This distinction is essential. What we find troubling in such encounters is not excess vitality but insufficient inhibition. The animals have not learned which expressions are appropriate to which contexts, nor how to suppress impulses that, while internally motivated, impose external costs. The analogy to aesthetic judgment is not metaphorical but structural.

\section{Taste as Negative Capacity}

Taste, in this framework, is best understood not as a faculty of appreciation but as a negative capacity: the learned ability to refrain from expression. To have taste is not merely to know what could be said, shown, or emphasized, but to know what should remain implicit, understated, or entirely absent. This is why technical proficiency alone is insufficient to guarantee aesthetic success. One may possess the skill to render detail, simulate texture, or encode symbolism, and yet fail aesthetically by exercising those skills indiscriminately.

Gaudiness arises when expressive systems behave like untrained animals: when they respond to every stimulus, amplify every affordance, and externalize every internally available feature. The resulting output may be energetic, intricate, and even impressive in isolation, but it lacks compositional discipline. It does not respect the interpretive limits of its audience, nor the contextual boundaries that make expression meaningful rather than noisy.

This insight allows us to reframe longstanding aesthetic disputes. The problem with certain caricatures, cartoons, or anime styles is not that they are stylized, artificial, or exaggerated. Stylization itself is a form of restraint. The problem is that exaggeration is often applied uniformly, without hierarchy or modulation. Every feature shouts. Nothing listens. Similarly, diagrams that attempt to include every relevant variable, annotation, or visual cue often fail not because they are inaccurate, but because they overwhelm the perceptual channel through which understanding must pass.

In such cases, expressiveness is confused with communication. The system behaves as though inclusion were equivalent to clarity, when in fact clarity often depends on omission.

\section{Outline of the Argument}

The remainder of this essay develops this account across several domains. In the next section, I examine historical portraiture to show how aesthetic power often derives from subtle deviations within otherwise constrained forms, rather than from symbolic accumulation. I then turn to contemporary information design and AI-generated imagery, arguing that their failures repeat long-standing errors under new technological conditions. Finally, I analyze AI-generated audio narration as a case study in expressive competence without social grounding, where fluency and insight coexist with errors no human speaker would naturally commit.

Across these cases, the same pattern recurs. Gaudiness is not the enemy of sophistication. It is sophistication without manners. It is expression untempered by judgment, and power exercised without restraint.

\section{Restraint and Suggestion in Historical Portraiture}

A useful counterpoint to contemporary accusations of aesthetic excess can be found in the history of portraiture. The most enduring portraits in Western art are rarely those that attempt exhaustive representation. Their power lies not in maximal detail, but in disciplined suggestion. Anatomical correctness, while often present, is subordinate to subtler compositional decisions: the modulation of line weight, the control of tonal gradients, the direction and softness of light, and the deliberate shaping of facial expression. These elements function not as decorations, but as constraints through which character becomes legible.

In many successful portraits, the sitter’s identity is not asserted through an inventory of attributes but inferred through controlled deviation. A slight asymmetry in the mouth may suggest irony or reserve. A tension around the eyes may indicate vigilance, fatigue, or authority. The posture of the body, even when minimally articulated, can communicate occupation, class, or temperament without recourse to explicit symbols. What matters is not that these signals are unmistakable, but that they are plausibly interpretable. They invite the viewer into an act of inference rather than overwhelming them with assertion.

This stands in contrast to iconographic traditions that rely heavily on symbolic accumulation. Medieval depictions of saints, for instance, often employ a dense visual vocabulary of objects, colors, wounds, and gestures, each encoding a specific narrative or theological reference. Such images are not primarily aesthetic in the modern sense; they are mnemonic and didactic. Their effectiveness depends on the viewer’s prior knowledge of the symbolic code. Without that code, the image collapses into clutter.

Even within this tradition, the limitations of excessive symbolism are apparent. Consider the frequent depiction of John the Baptist holding his severed head. The image is intended to compress biography into a single visual token, yet it introduces problems that are both aesthetic and conceptual. Anatomically implausible and visually grotesque, the motif risks reducing a complex figure to the spectacle of his death. The symbol overwhelms the subject. The image ceases to suggest character and instead insists upon a single, lurid fact.

What distinguishes later portrait traditions is not a rejection of symbolism but a rebalancing of it. Recognizable elements such as clothing, color, or modest props persist, but they are subordinated to formal restraint. The symbol becomes an accent rather than a scaffold. Meaning emerges from the interaction between what is shown and what is withheld.

\section{Correctness Versus Judgment}

The historical examples illustrate an important distinction between correctness and judgment. A portrait can be anatomically correct and aesthetically disastrous. Conversely, a portrait can deviate from strict realism and achieve profound expressive clarity. The difference lies not in fidelity to external standards but in the internal coherence of the expressive system.

This distinction is frequently lost in contemporary discussions of AI-generated imagery. Critics often conflate technical competence with aesthetic legitimacy, assuming that improvements in realism, resolution, or anatomical accuracy will automatically yield better art. Yet history suggests the opposite. Technical mastery expands the space of possible expression, but it simultaneously increases the burden of selection. The more one can render, the more one must choose not to render.

Gaudiness arises when this burden is evaded. When every available feature is exercised, when every stylistic affordance is activated, the result is not richness but noise. The image behaves like an untrained system, responding to every internal impulse without regard for external legibility. The failure is not one of capability but of inhibition.

This helps explain why certain technically accomplished illustrations and diagrams remain difficult to parse. They may contain all relevant information, rendered with precision and care, yet fail as communicative artifacts. The problem is not that they include too much information per se, but that they lack a hierarchy of importance. Without such hierarchy, the viewer is forced to perform the work of selection that the designer has abdicated.

In this sense, aesthetic judgment is inseparable from ethical responsibility. To present information, images, or expressions to others is to impose a cognitive load. Taste consists in managing that load responsibly. Gaudiness is what happens when the creator externalizes their own excitement and curiosity without translating it into a form that others can comfortably inhabit.

\section{From Symbolic Excess to Algorithmic Excess}

The transition from medieval iconography to contemporary algorithmic imagery may appear radical, but the underlying failure mode is continuous. Both rely on the accumulation of signifiers without sufficient modulation. Where the medieval image piles symbols, the modern AI-generated image piles stylistic effects: exaggerated lighting, hyper-saturated color, improbable textures, and expressive distortions layered without restraint.

Importantly, this is not an indictment of stylization itself. Stylization, when disciplined, is a powerful form of abstraction. The problem arises when stylization is treated as a license rather than a constraint. In many AI-generated portraits, every stylistic dimension is pushed toward maximal expressiveness simultaneously. The face glows, the eyes shimmer, the background erupts with symbolic clutter, and the emotional tone is amplified beyond any plausible social context.

Such images resemble the barking dog at the window. They are eager, insistent, and poorly calibrated. They demand attention rather than earning it. The viewer is not invited into interpretation but assaulted by assertion.

This pattern also explains why caricature and certain cartoon traditions provoke discomfort rather than delight. Exaggeration is not inherently grotesque; it becomes grotesque when it lacks internal limits. When every feature is distorted, distortion itself loses meaning. The image no longer points beyond itself; it collapses into self-parody.

\section{The Preparatory Role of Restraint}

What these cases collectively demonstrate is that restraint is not a secondary refinement applied after expressive capacity is achieved. It is a foundational condition of intelligibility. Without restraint, expression does not merely become excessive; it becomes socially incoherent.

This insight prepares the ground for the analysis of AI-generated audio narration, where expressive fluency and conceptual understanding coexist with failures of situational awareness. In that domain, as in visual art, the problem is not artificiality but a lack of lived calibration. The system knows how to speak, but not when not to speak, nor which expressions presuppose experiences it cannot legitimately claim.

The discussion thus far has treated gaudiness primarily as a visual phenomenon, rooted in symbolic overload and insufficient modulation. The same structural failure, however, appears in non-visual media. When expressive systems acquire fluency without social grounding, they reproduce in speech the same errors that untrained animals and undisciplined artists reproduce in behavior and image-making. AI-generated audio summaries and podcast-style narrations provide a particularly clear case, as their near-human competence makes the absence of contextual restraint especially salient.

\section{Fluency Without Grounding: AI-Generated Audio and the Performance of Familiarity}

The failures of restraint described thus far become particularly visible in the domain of AI-generated audio narration, especially in formats designed to resemble human podcasts or reflective essays. These systems often exhibit remarkable linguistic competence. They accurately summarize complex material, identify salient arguments, and reproduce patterns of emphasis and cadence that closely resemble human speech. Their pronunciation and prosody can be subtle, expressive, and in many cases more carefully modulated than that of casual human speakers.

Yet it is precisely in this near-success that the aesthetic failure emerges. The problem is not that these narrations are artificial, nor that they lack intelligence. It is that they perform social and temporal claims that presuppose experiences they do not possess. The result is not confusion but a distinct form of tackiness: an expressive overreach that violates shared expectations about who is entitled to say what.

A recurring example occurs when an AI-generated overview introduces an essay written days earlier by claiming that its ideas have been “stuck in my head all week” or that the material “kept me up all night.” These phrases are idiomatic signals of intellectual or emotional engagement. They are not factual claims in the narrow sense, but they do rely on an implicit temporal narrative. A human speaker uses them to situate themselves within a lived sequence of reflection, boredom, return, and persistence. When such phrases are deployed without regard to temporal plausibility, the performance collapses.

What makes these errors striking is that they are not mistakes of comprehension. The system understands the function of the phrase. It knows that such expressions convey depth of engagement. What it lacks is the capacity to withhold expressions that require biographical continuity. The failure is not semantic but situational. The system does not know which idioms are socially licensed and which must be earned through lived duration.

This is why such outputs feel tacky rather than merely wrong. A factual error can be corrected. A stylistic infelicity can be forgiven. But a misplaced claim of familiarity produces discomfort because it mimics intimacy without justification. It is the verbal analogue of standing too close to a stranger or using a nickname that has not been offered.

\section{Name, Pronunciation, and the Limits of Archive-Based Speech}

A related class of failures concerns proper names, particularly those that fall outside well-represented phonetic archives. The persistent mispronunciation of my author name, Flyxion, is a minor but illustrative case. Variants such as “Flection,” “Fliction,” “Affliction,” “Flishannon,” or “Flycannon” are understandable attempts to map an unfamiliar orthographic form onto existing phonological patterns. A human encountering the name for the first time might hesitate or ask for clarification.

What distinguishes the AI-generated errors is not their frequency but their confidence. The system does not mark uncertainty. It does not pause, hedge, or request correction. It selects a plausible-sounding pronunciation and proceeds as though familiarity had been established. This again is not a technical failure but a social one. Human speakers routinely signal uncertainty when encountering unfamiliar names, precisely because names carry personal and social weight. To mispronounce a name while signaling confidence is not merely inaccurate; it is impolite.

The issue becomes more consequential when applied to culturally and historically significant names. In discussing my essay on a screenplay centered on the Flower Wars, AI-generated narration has repeatedly mispronounced “Mexica,” rendering it as “Mecksica” rather than the historically grounded pronunciation closer to “Mesheeca.” This is not a trivial accent variation. The Mexica possessed a phoneme equivalent to the English “sh,” which Spanish lacked. Early Spanish scribes approximated this sound using the letter \emph{x}, then a relatively flexible grapheme. Over time, Spanish phonology shifted, but the orthography persisted, producing the modern pronunciation of \emph{México} with an aspirated sound.

English speakers, encountering the letter \emph{x} without this historical context, often default to a hard consonant. The AI system reproduces this default without hesitation, despite having access to the relevant linguistic history. The failure is not informational; it is one of salience and restraint. The system does not recognize that this is a case where defaulting is inappropriate, where correction matters, and where mispronunciation reproduces a long-standing pattern of cultural flattening.

Once again, the discomfort produced is not rooted in error alone. It arises from expressive overconfidence. The system asserts familiarity where humility would be appropriate. It performs authority without accountability.

\section{Tackiness as Performed Intimacy}

Across these examples, a common structure emerges. Gaudiness and tackiness are not simply aesthetic defects; they are failures of social positioning. They occur when expressive systems adopt tones, gestures, or claims that exceed their legitimate standing. In visual art, this appears as symbolic overload. In diagrams, as informational clutter. In AI-generated audio, as the performance of intimacy, reflection, or cultural fluency without the lived grounding that makes such performances credible.

This explains why such outputs often provoke a reaction that is difficult to articulate yet immediately felt. They are not obviously incompetent. On the contrary, they are often impressively capable. What they lack is an internalized sense of when to step back. They do not know how to greet without lunging, how to indicate interest without overwhelming, or how to signal understanding without claiming experience.

Taste, in this sense, is inseparable from social calibration. It is the ability to align expression with position, capacity, and context. Gaudiness is what results when that alignment fails.

\section{The Ethical Dimension of Restraint}

The argument thus far suggests that restraint is not merely an aesthetic virtue but an ethical one. To exercise restraint is to acknowledge the presence of others, their limits, and their interpretive labor. It is to take responsibility for how one’s expressions will be received, not merely for their internal coherence.

This ethical dimension becomes especially salient in systems designed to scale expression beyond individual human limits. When expressive capacity is amplified without a corresponding amplification of restraint, the burden of interpretation is shifted onto audiences. They must filter, discount, and reinterpret outputs that could have been moderated at the source.

Gaudiness, then, is not harmless exuberance. It is a form of cost externalization. It demands attention, correction, or endurance from others. In this sense, it resembles other forms of unregulated behavior that are tolerated only because their harms are diffuse.

Having examined behavioral, visual, and auditory cases, the underlying structure of tackiness should now be clear. It is not the presence of expression but the absence of inhibition. It is not artificiality but unearned familiarity. It is not excess energy but insufficient judgment.

\section{Abstraction, Error, and the Continuous Renegotiation of Meaning}

It would be a mistake to interpret the preceding analysis as an argument against abstraction, simplification, or mediated expression. All communication, whether visual, textual, or auditory, depends on abstraction. To speak at all is to omit detail; to draw is to suppress dimensions; to summarize is to distort. The presence of error, oversimplification, or approximation is therefore not a defect introduced by computation, but a constitutive feature of meaning-making itself.

Human expression is never exact. Any attempt to describe a topic, especially one of complexity or novelty, will contain minor inaccuracies, imprecisions, slips of pronunciation, or stylistic infelicities. These are ordinarily tolerated, and often go unnoticed, so long as they do not interfere with comprehension. Indeed, a certain level of noise is expected. Communication functions not because it eliminates error, but because it allows interlocutors to correct, contextualize, and negotiate meaning in real time.

What distinguishes acceptable error from jarring error is not its existence but its impact on interpretability. A mispronounced word that remains recognizable is rarely distracting. A spelling mistake that preserves phonetic structure is quickly repaired by the reader. By contrast, an error that interrupts parsing, misdirects inference, or forces the audience to stop and reconstruct intent imposes a cognitive cost. It is at this point that approximation ceases to be invisible and becomes an aesthetic problem.

Crucially, the kinds of errors that appear are shaped by the medium through which expression is produced. Writing on a mechanical typewriter encourages a specific error profile: adjacent keys may be struck accidentally, or the metal bars corresponding to letters may jam, producing visible collisions or omissions. These errors are material, mechanical, and easily attributed. Their origin is legible, and therefore rarely taken as evidence of conceptual failure.

Digital interfaces introduce different distortions. Autocorrect systems, predictive text, and gesture-based input methods impose their own priors on expression. On a trace-based keyboard, for instance, the word \emph{fluxion} may be repeatedly altered into \emph{Flyxion}, not because the former is incorrect—it is a legitimate term with historical pedigree—but because the latter has been used more frequently and more recently. The system optimizes for local statistical salience rather than global lexical correctness. The resulting error is not random but adaptive, and its source lies in usage patterns rather than intent.

By contrast, when touch typing on a QWERTY keyboard, such an error may never occur, though others may arise instead. A prolonged press of the shift key may yield an unintended capital letter, producing \emph{FLyxion} rather than \emph{Flyxion}. The error is different in kind, arising from motor timing rather than lexical substitution. In each case, the mistake is not evidence of misunderstanding. It is a byproduct of the interface through which meaning is externalized.

These observations generalize. Speech errors differ from writing errors; handwriting differs from typesetting; live conversation differs from edited prose. Fatigue, distraction, and familiarity further modulate error profiles. A speaker who slept poorly may stumble over articulation. A writer working quickly may omit articles or conjunctions. None of these deviations are inherently problematic. They become problematic only when they disrupt the listener’s or reader’s ability to reconstruct intent.

This is why it is misleading to single out computer-generated abstraction as uniquely prone to aesthetic failure. The issue is not that machines abstract, but that abstraction always carries risk. Meaning does not reside fully in any single utterance or artifact. It emerges through continuous renegotiation between producer, medium, and audience. Each new expressive layer introduces new opportunities for both clarity and misalignment.

What distinguishes tasteful abstraction from tacky abstraction is not the absence of error, but the presence of corrective affordances. Human communicators instinctively hedge, pause, revise, or solicit feedback when uncertainty arises. They mark approximation as approximation. They adjust in response to confusion. Systems that lack these feedback-sensitive brakes tend to press forward with confident approximations, even when those approximations strain credibility or comprehension.

Thus, the problem identified throughout this essay is not abstraction per se, nor simplification, nor mediation by machines. It is the failure to recognize that meaning is provisional. Expression must remain negotiable. When abstraction is treated as final rather than tentative, when approximation is presented as authority, and when errors are delivered without cues for repair, the result is not merely inaccuracy but aesthetic discomfort.

Gaudiness, in this expanded sense, is what happens when abstraction forgets that it is abstracting. It is expression that mistakes compression for completion. Taste, by contrast, preserves room for adjustment. It acknowledges that communication is a cooperative process, not a unilateral broadcast, and that meaning survives not through precision alone but through the shared ability to recover from error.

\section{Titles, Separators, and the Infrastructure of Expression}

The question of restraint extends beyond images, speech, and abstraction to the infrastructural layer of expression itself. Titles, filenames, and metadata fields are not neutral containers for meaning; they are interfaces governed by technical conventions that shape how work is stored, retrieved, and encountered. Decisions that appear purely stylistic at the level of prose often have unintended consequences once they interact with software systems, file hierarchies, and display constraints.

Commas provide a simple example. In ordinary writing, the comma is a tool of nuance, allowing clauses to be balanced and distinctions to be refined. In programming and data representation, however, the comma functions primarily as a separator. It signals the boundary between arguments, fields, or elements in a list. When a title containing commas is repurposed as a filename, URL slug, database key, or command-line argument, that semantic subtlety is lost. The punctuation is reinterpreted syntactically, and what was meant to bind ideas together instead fractures them.

This is not a theoretical concern. Many systems implicitly tokenize strings on commas. Others escape them inconsistently. Some treat them as reserved characters; others allow them but display them unpredictably. The result is a class of errors that are not semantic in origin but infrastructural in effect. A title that reads clearly on the page may behave poorly when passed through a toolchain, producing broken links, malformed references, or silent truncation.

Length introduces a related problem. While modern publishing platforms technically permit long titles, the surrounding infrastructure often does not honor that freedom. Titles are routinely truncated in tables of contents, browser tabs, PDF viewers, citation managers, search results, and mobile interfaces. What survives these truncations is rarely the carefully balanced structure of the full title, but an arbitrary prefix. Excessive elaboration at the title level therefore risks misrepresenting the work at precisely the moments when identification matters most.

Here again, the issue is not expressiveness but calibration. A title is not merely a sentence; it is an address. It must function across multiple layers of representation, from human reading to machine parsing to visual display. To load it with subordinate clauses, punctuation-dependent nuance, or extended rhetorical flourish is to ignore the constraints of the systems through which it must travel.

This helps explain why shorter titles with explicit subtitles often outperform longer, unified constructions. The main title carries the durable identifier, while the subtitle absorbs explanatory detail. The structure mirrors good design elsewhere: a stable handle paired with optional elaboration. Restraint at the infrastructural level preserves meaning by ensuring that what remains visible under constraint is still accurate.

The relevance of this observation to the broader argument should now be apparent. Gaudiness does not only occur at the level of image or sound; it also appears in metadata. A title that attempts to say everything at once may satisfy the author’s desire for precision while failing its primary function as a point of reference. Taste, once again, consists in knowing which expressive ambitions to defer.

In this sense, the decision to avoid commas or excessive length in a title is not an act of concession to machinery, but an acknowledgment of context. It reflects an understanding that meaning must survive translation across systems, and that expressive discipline includes anticipating the environments in which an artifact will circulate.

\section{Conclusion: Taste as Disciplined Negation}

The cases examined throughout this essay—untrained animal behavior, symbolic overload in historical imagery, contemporary visual excess, and AI-generated audio narration—share a single structural failure. In each instance, expressive capacity outpaces the mechanisms that regulate its deployment. The result is not error in the narrow sense, nor incompetence, but a breakdown in calibration. Gaudiness and tackiness arise when systems speak, show, or signal without sufficient regard for context, hierarchy, or social standing.

This reframing allows us to move beyond the common temptation to locate aesthetic failure in technology itself. Artificial intelligence did not invent bad taste, nor did mass production, nor even modernity. Each new expressive medium merely exposes, at greater scale and clarity, a problem that has always accompanied expressive power. Whenever the cost of expression drops faster than the cost of judgment, tackiness proliferates.

Taste, on this account, is not a mysterious faculty or an innate sensitivity. It is a learned discipline of omission. To have taste is to know not only what one can express, but which expressions impose unnecessary burdens on others. It is the capacity to suppress signals that are internally motivated yet externally disruptive. This is why restraint is not the enemy of expressiveness but its enabling condition. Without restraint, expression collapses into noise.

The analogy with social behavior is instructive. The barking dog is not objectionable because it is alive, excited, or friendly. It is objectionable because it has not learned how to greet without overwhelming. The same holds for expressive systems. A portrait that insists on every symbolic association, a diagram that refuses to rank its variables, or an audio narration that claims lived reflection it cannot possess all exhibit the same defect. They behave as though all available signals deserve equal realization.

What distinguishes sophistication from gaudiness is not complexity but hierarchy. Sophisticated systems know which features must dominate, which must recede, and which must remain entirely implicit. They understand that meaning often emerges from tension between presence and absence. Gaudiness erases that tension by flattening expressive priority.

This perspective also clarifies why certain aesthetic failures provoke irritation rather than mere indifference. Gaudiness is not passive. It demands attention. It insists on being processed. In doing so, it externalizes the work of restraint onto its audience. Viewers, listeners, or readers are forced to filter what should have been filtered at the point of production. The offense lies not in exuberance but in the conscription of others’ attention.

In the context of artificial intelligence, this analysis carries an important implication. The challenge facing generative systems is not primarily one of improving realism, coherence, or stylistic variety. These capabilities already exist in abundance. The more difficult task is to encode norms of restraint: when to hedge, when to remain silent, when to signal uncertainty, and when to forgo expressive flourishes that presuppose experiences the system cannot legitimately claim.

This is not a purely technical problem. It is a cultural and ethical one. Restraint cannot be learned from archives alone, because archives disproportionately record assertive speech. They preserve declarations, performances, and artifacts, but rarely capture the countless moments of omission that constitute good judgment. Taste is transmitted less through explicit rules than through prolonged participation in social feedback loops where miscalibration carries consequences.

The persistence of tackiness in new media, then, should not surprise us. What is surprising is how clearly it reveals the nature of taste itself. Taste is not about knowing what to add. It is about knowing what to leave out. It is not about refinement as ornament, but refinement as refusal.

In this sense, gaudiness is not the opposite of intelligence, nor even of creativity. It is intelligence without humility, creativity without inhibition, and expression without manners. The cultivation of taste—whether in humans, institutions, or machines—begins not with amplification, but with the disciplined practice of saying less.



\end{document}
